\NativeSlot{47}{55}{%
\ExamGroup{三、计算。{\normalfont\kaishu\fontsize{10.8bp}{14bp}\selectfont （共26分）}}
\noindent 1.\quad 直接写得数。（8分）\par
\noindent%
% p031_s01
\NativeOralCell{28+172=}%
% p031_s02
\NativeOralCell{136-45-55=}%
% p031_s03
\NativeOralCell{200\div25\div2=}%
% p031_s04
\NativeOralCell{160\div2\times5=}%
\par
\noindent%
% p031_s05
\NativeOralCell{101\times12=}%
% p031_s06
\NativeOralCell{139+65-39=}%
% p031_s07
\NativeOralCell{37\times99+37=}%
% p031_s08
\NativeOralCell{8\times5\div8\times5=}%
\par

}
\NativeSlot{146}{72}{%
\noindent \textbf{2.}\quad 用简便方法计算下面各题。（18分）\par\smallskip
\par\noindent%
\ExamTile{.323}{%
% p031_s09

\noindent （1）\(138+293+62+107\)\par


}\hfill%
\ExamTile{.323}{%
% p031_s10

\noindent （2）\(574-189-211\)\par


}\hfill%
\ExamTile{.323}{%
% p031_s11

\noindent （3）\((25+7)\times8\)\par


}\par
}
\NativeSlot{297}{63}{%
\par\noindent%
\ExamTile{.323}{%
% p031_s12

\noindent （4）\(250\times57\times40\)\par


}\hfill%
\ExamTile{.323}{%
% p031_s13

\noindent （5）\(720\div16\div5\)\par


}\hfill%
\ExamTile{.323}{%
% p031_s14

\noindent （6）\(45\times102\)\par


}\par
}
\NativeSlot{421}{99}{%
\ExamGroup{四、动手操作。{\normalfont\kaishu\fontsize{10.8bp}{14bp}\selectfont （共6分）}}

% p031_s15

\ExamQuestion{操作}{7}%
请你回忆学习“乘法分配律”时的探究过程，举例说明“\(12\times9+8\times9=(12+8)\times9\)”，可以用画图、列表、文字叙述等表达你的思考过程。（6分）\par
}
\NativeSlot{597}{107}{%
\ExamGroup{五、解决问题。{\normalfont\kaishu\fontsize{10.8bp}{14bp}\selectfont （共27分）}}

% p031_s16

\ExamQuestion{1.}{7}%
小军看一本故事书，第一天看了178页，第二天看了63页，还剩下122页没看，这本故事书一共有多少页？（5分）\par
}
